\documentclass{zhvt-classic}

% 使用本地源流明体（GenRyuMin2 TW）按文件路径加载，不依赖系统字体名
% 版框（内框／界栏／书口侧栏／鱼尾）改用与版心印章同一印泥棕色：
%   印章图片 yshs.png 实色 = RGB(135,68,52) = HTML 874434
\usepackage{xcolor}
\usepackage{tikz}
% sealbrown 双注册：xcolor 供 TikZ/pgf，L3 供 cls 版框
\definecolor{sealbrown}{HTML}{874434}
\ExplSyntaxOn
\color_set:nnn { sealbrown } { HTML } { 874434 }
\ExplSyntaxOff

\zhvtset{
  grid color       = sealbrown,
  main font        = GenRyuMin2TW-SB,
  main font path   = C:/Users/Administrator/texmf/fonts/opentype/public/GenRyuMin2TW/,
jiazhu font      = GenRyuMin2TW-M,
  jiazhu font path = C:/Users/Administrator/texmf/fonts/opentype/public/GenRyuMin2TW/,
  % 页码框容纳 4 字：正常页「九十七」+前置空 1 格=4 行恰好；
  % 100+ 页「一百零一」共 5 行超出，末字「一/二/三」被框裁掉（2026-10-04 实测），故放宽到 5
  shukou page number chars = 5,
  % 书口书名框高从默认 6 字减到 4 字，把省出的 2 字高拨给 chapter 框，
  % 使 6 字书名+卷名放进 9 字高的 chapter 框后不再压线（2026-10-05）
  shukou title chars = 4,
}

\title{史记}[燕召公世家]
\author{汉·司马迁}
\maker{闲叟漫记·重製}


% ===== 索隐述赞样式（2026-11-08 四次定稿 · 愚大人亲手改定参数）=====
% ⚠ 竖排前提：自定义字体家族必须带 RawFeature={vertical:+vert}；
%   simhei.ttf 加 +vhal 会整体加载失败静默回退，只写 +vert；simfang 同理。
\newCJKfontfamily{\shuzanfont}[Path=C:/Windows/Fonts/,RawFeature={vertical:+vert}]{simhei.ttf}
\newCJKfontfamily{\shuzanbody}[Path=C:/Windows/Fonts/,RawFeature={vertical:+vert}]{simfang.ttf}
% 题头：152x30pt 圆角 6pt 棕底白字；baseline (0,2) 使整框落入界行格子不压版框线；
% 白字节点 x 不补偿、y +8.2pt（像素级墨迹实测 dx/dy≈0，L/R 留白约 20pt）。
\newcommand{\shuzanhead}[1]{%
  \begin{tikzpicture}[baseline={(0,2)}]
    \fill[sealbrown, rounded corners=6pt]
      (-76pt,-15pt) rectangle (76pt,15pt);
    \node[text=white, inner sep=0pt, anchor=center] at (0pt,8.2pt)
      {\shuzanfont\fontsize{23}{23}\selectfont #1};
  \end{tikzpicture}%
}
% ★ 修正书口卷名位置：chapter 框起始位置改回 cls 默认值，避免全局影响。
%   只在每个卷首页通过 \juan 命令里的空命令微调该页位置（安全，不会跨页污染）。
\ExplSyntaxOn
\AtBeginDocument{
  \dim_set:Nn \l__zhvt_shukou_chapter_title_start_dim
    { \l__zhvt_shukou_title_height_dim + \l__zhvt_font_size_dim * 2 }
  \dim_set:Nn \l__zhvt_shukou_page_number_start_dim
    { \l__zhvt_shukou_chapter_title_start_dim + \l__zhvt_shukou_chapter_title_height_dim }
}
\ExplSyntaxOff

\begin{document}

% ★★ 版心（书口）排布覆写：把书名放到鱼尾之下（2026-10-04）
%
%   cls 默认版心自上而下：书名（title 静态框，在鱼尾之上）→ 鱼尾 → 篇名（chapter 动态框）→ 页码
%   本模板改为：          鱼尾（位置不动）→ 书名 → 篇名 → 页码分割线 → 页码
%
%   手法：把 cls 放在 title 静态框里的书名，改放进 chapter 动态框（鱼尾之下）；
%         title 静态框留空。书名取 \l__zhvt_title_str，即 \title{...} 的第一参数，
%         所以换书时只改 \title，不必在此重复写书名。
%
%   ⚠ 三条必须遵守的写法（2026-10-04 实测）：
%     ① 汉字必须用 \hbox{...} 包住。flowfram 会把框内容装进一个宽 14.4mm（= 一个 line_sep
%        = 书口栏宽）的纵向 minipage 并按该宽度折行，而整页靠 PDF /Rotate 90 显示、
%        框内文字本身不变向 —— 不包 \hbox 时「围炉夜话」4 字会被逐字折成 4 行，
%        在版面上就成了横排 4 字并横向溢出版框。
%     ② 竖排下**一个 \hbox 就是一整列**。书名与篇名要同处一列，必须放进同一个 \hbox；
%        分成两个 \hbox 时，后一个会被当成新的一列、甩到隔壁栏去
%        （实测「序」跑到书口栏外侧 14.4mm 的正文列上）。中间的 \quad 在框内是横向空白，
%        显示出来正好是书名与篇名之间的纵向间隔。
%     ③ ★ 内容盒要写成 **与框等宽的定宽盒** `\hbox to \l__zhvt_shukou_chapter_title_height_dim{…\hfil}`。
%        实测：chapter 框是把内容盒沿纵向**居中**摆放的，内容宽度一变，书名就跟着漂：
%        自然宽度的 \hbox（书名 38.4mm）落在 x=157.2mm，加篇名后（57.6mm）落在 x=152.4mm，
%        目录页（无篇名）与正文页的书名就差半格。做成等宽盒后，无论有没有篇名，
%        内容固定从框顶起排，实测两页均落在 x=142.8mm，不再漂。
\ExplSyntaxOn
% 把书口 chapter 框（书名+篇名）起始位置下移 1 字高，避免 4 字卷名顶到鱼尾
% ⚠ cls 的 \__zhvt_calculate_dimensions: 在 \AtBeginDocument 里执行，会重算该 dim；
%   必须把覆写挂到 begin document 之后，否则会被二次覆盖（2026-10-05 实测）。
% (旧方案已移到导言区，直接修改 chapter 框起始位置)
% 卷标题命令 #1=书口篇名(鱼尾) #2=正文卷名（含卷序，卷名内空格用 ~）
%
% ★★ 目录页两项修订（2026-10-05，愚大人指出）★★
%   ① 页码起点统一对齐：
%      原状 —— 目录项写成 \par\hspace{4em}卷名\hspace{4em}页码，页码起点 = 卷名实宽 + 4em。
%      卷名实宽随字数变：「卷一~梁惠王上」6 字 +~、「卷七~离娄上」5 字 +~，
%      差 1 个字位 → 卷七~卷十 4 条的页码整体比其余 10 条高 27.21pt（一个字），看着上下不齐。
%      改法 —— 卷名补全角空格（宽 = 1 字 = 27.215pt）补足到 **7 字符位**再写进目录，
%      14 条页码起点遂全部落在同一条竖线上（起点对齐；末字随页码字数自然参差）。
%      ⚠ 只对目录里的副本 \l_tmpc_tl 补，正文卷名仍用 \l_tmpa_tl 原样 —— 正文像素不受影响。
%   ② 卷名整体上移半字位：前导 \hspace 由 4em 改 3.5em。
%      实测原状目录块在版框内偏下（天头留白 139.6pt / 地脚留白 110.7pt），
%      上移 0.5em = 13.6pt 后上下基本均分，视觉居中。
\NewDocumentCommand { \juan } { mm } {
	\clearpage
	\tl_set:Nn \l__zhvt_chapter_title_tl { #1 }
	\tl_set:Nn \l_tmpa_tl { #2 }
	\tl_set_eq:NN \l_tmpc_tl \l_tmpa_tl
	\int_step_inline:nn { 7 - \str_count:n { #2 } }
	  { \tl_put_right:Nn \l_tmpc_tl { 　 } }
	\tl_set:Nx \l_tmpb_tl { \zhnum { page } }
	\tl_put_left:Nn \l_tmpb_tl { \hspace { 4em } }
	\tl_put_left:Nn \l_tmpb_tl { \l_tmpc_tl }
	\tl_put_left:Nn \l_tmpb_tl { \par \hspace { 3.5em } }
	\addtocontents { toc } { \l_tmpb_tl }
	\skip_horizontal:n { \l__zhvt_font_size_dim * 4 }
	\l_tmpa_tl
	\par
	\thispagestyle { empty }
}
\cs_set_protected:Npn \__zhvt_set_shukou_content: {
	% ★ 鱼尾上方留白（title 静态框）嵌入闲叟钤印：
	%   整页 /Rotate 90，故此处把图用 \rotatebox{90} 转正为竖章；
	%   尺寸卡在版心条幅内（书内约 7mm 宽 × 33.8mm 高，留余量不压鱼尾）。
	\setstaticcontents * { title } {
		\vbox_to_ht:nn { \l__zhvt_baseline_skip_dim } {
			\vskip 14.2pt
			\hbox_to_wd:nn { \l__zhvt_shukou_title_height_dim } {
				\hfil
				\rotatebox [ origin = c ] { 90 } { \includegraphics [ width = 7mm , height = 33.8mm ] { yshs.png } }
				\hfil
			}
		}
	}
% ★ 版心（书口）统一改法（2026-11-16 愚大人定）：18pt＋框内居中＋书口文字同版框色
%   ① 原 27.2pt 时 7 字长篇名（如「廉颇蔺相如列传」）末字溢出 chapter 框、压住页码；
%      改 \fontsize{18pt}{18pt}，前后 \hfil 令字串在框内居中（不随篇名字数漂移）。
%   ② \vbox{\vskip 8.2pt} 定量下移，抵消缩字号后栏宽方向 8.2pt 上偏（像素实测 y 中心 543.3 = 原版）。
%      ⚠ \raisebox / \vspace* 在 flowfram 框内无效；换字号须同步换补偿：16pt→10.1、19pt→7.4。
%   ③ chapter 与 page~number 均置 \color{sealbrown}（HTML 874434），与版框／界行同色；正文仍黑。
%   ⚠ 页码字位、钤印、鱼尾、界行、正文一律未动；五部合订本同此口径。
	\setdynamiccontents * { chapter }
	{ \color{sealbrown} \vbox { \vskip 8.2pt \hbox to \l__zhvt_shukou_chapter_title_height_dim
		{ \hfil \fontsize{18pt}{18pt}\selectfont \l__zhvt_title_str \quad \l__zhvt_chapter_title_tl \hfil } } }
	\setdynamiccontents * { page~number } { \color{sealbrown} \quad \zhnum { page } }
}
\__zhvt_set_shukou_content:
\ExplSyntaxOff

\maketitle

% 单篇，无多卷目录，不渲染目录页

\juan{燕召公世家}{卷三十四~燕召公世家第四}
\hspace*{2em}召公奭与周同姓。姓姬氏。周武王之灭纣。封召公於北燕。\\
\hspace*{2em}其在成王时。召王为三公。自陕以西。召公主之。自陕以东。周公主之。成王既幼。周公摄政。当国践祚。召公疑之。作君奭。君奭不说周公。周公乃称﹁汤时有伊尹。假于皇天。在太戊时。则有若伊陟。臣扈。假于上帝。巫咸治王家。在祖乙时。则有若巫贤。在武丁时。则有若甘般。率维兹有陈。保乂有殷﹂。於是召公乃说。\\
\hspace*{2em}召公之治西方。甚得兆民和。召公巡行乡邑。有棠树。决狱政事其下。自侯伯至庶人各得其所。无失职者。召公卒。而民人思召公之政。怀棠树不敢伐。哥咏之。作甘棠之诗。\\
\hspace*{2em}自召公已下九世至惠侯。燕惠侯当周厉王奔彘。共和之时。\\
\hspace*{2em}惠侯卒。子厘侯立。是岁。周宣王初即位。厘侯二十一年。郑桓公初封于郑。三十六年。厘侯卒。子顷侯立。\\
\hspace*{2em}顷侯二十年。周幽王淫乱。为犬戎所弑。秦始列为诸侯。\\
\hspace*{2em}二十四年。顷侯卒。子哀侯立。哀侯二年卒。子郑侯立。郑侯三十六年卒。子缪侯立。\\
\hspace*{2em}缪侯七年。而鲁隐公元年也。十八年卒。子宣侯立。宣侯十三年卒。子桓侯立。桓侯七年卒。子庄公立。\\
\hspace*{2em}庄公十二年。齐桓公始霸。十六年。与宋。卫共伐周惠王。惠王出奔温。立惠王弟颓为周王。十七年。郑执燕仲父而内惠王于周。二十七年。山戎来侵我。齐桓公救燕。遂北伐山戎而还。燕君送齐桓公出境。桓公因割燕所至地予燕。使燕共贡天子。如成周时职。使燕复修召公之法。三十三年卒。子襄公立。\\
\hspace*{2em}襄公二十六年。晋文公为践土之会。称伯。三十一年。秦师败于殽。三十七年。秦穆公卒。四十年。襄公卒。桓公立。\\
\hspace*{2em}桓公十六年卒。宣公立。宣公十五年卒。昭公立。昭公十三年卒。武公立。是岁晋灭三郤大夫。\\
\hspace*{2em}武公十九年卒。文公立。文公六年卒。懿公立。懿公元年。齐崔杼弑其君庄公。四年卒。子惠公立。\\
\hspace*{2em}惠公元年。齐高止来奔。六年。惠公多宠姬。公欲去诸大夫而立宠姬宋。大夫共诛姬宋。惠公惧。奔齐。四年。齐高偃如晋。请共伐燕。入其君。晋平公许。与齐伐燕。入惠公。惠公至燕而死。燕立悼公。\\
\hspace*{2em}悼公七年卒。共公立。共公五年卒。平公立。晋公室卑。六卿始强大。平公十八年。吴王阖闾破楚入郢。十七年卒。简公立。简公十二年卒。献公立。晋赵鞅围范。中行于朝歌。献公十二年。齐田常弑其君简公。十四年。孔子卒。二十八年。献公卒。孝公立。\\
\hspace*{2em}孝公十二年。韩。魏。赵灭智伯。分其地。三晋强。\\
\hspace*{2em}十五年。孝公卒。成公立。成公十六年卒。愍公立。愍公三十一年卒。厘公立。是岁。三晋列为诸侯。\\
\hspace*{2em}厘公三十年。伐败齐于林营。厘公卒。桓公立。桓公十一年卒。文公立。是岁。秦献公卒。秦益强。\\
\hspace*{2em}文公十九年。齐威王卒。二十八年。苏秦始来见。说文公。文公予车马金帛以至赵。赵肃侯用之。因约六国。为从长。秦惠王以其女为燕太子妇。\\
\hspace*{2em}二十九年。文公卒。太子立。是为易王。\\
\hspace*{2em}易王初立。齐宣王因燕丧伐我。取十城。苏秦说齐。使复归燕十城。十年。燕君为王。苏秦与燕文公夫人私通。惧诛。乃说王使齐为反间。欲以乱齐。易王立十二年卒。子燕哙立。\\
\hspace*{2em}燕哙既立。齐人杀苏秦。苏秦之在燕。与其相子之为婚。而苏代与子之交。及苏秦死。而齐宣王复用苏代。燕哙三年。与楚。三晋攻秦。不胜而还。子之相燕。贵重。主断。苏代为齐使于燕。燕王问曰。﹁齐王傒如。﹂对曰。﹁必不霸。﹂燕王曰。﹁何也。﹂对曰。﹁不信其臣。﹂苏代欲以激燕王以尊子之也。于是燕王大信子之。子之因遗苏代百金。而听其所使。\\
\hspace*{2em}鹿毛寿谓燕王。﹁不如以国让相子之。人之谓尧贤者。以其让天下于许由。许由不受。有让天下之名而实不失天下。今王以国让于子之。子之必不敢受。是王与尧同行也。﹂燕王因属国于子之。子之大重。或曰。﹁禹荐益。已而以启人为吏。及老。而以启人为不足任乎天下。传之于益。已而启与交党攻益。夺之。天下谓禹名传天下于益。已而实令启自取之。今王言属国于子之。而吏无非太子人者。是名属子之而实太子用事也。﹂王因收印自三百石吏已上而效之子之。子之南面行王事。而哙老不听政。顾为臣。国事皆决于子之。\\
\hspace*{2em}三年。国大乱。百姓恫恐。将军市被与太子平谋。将攻子之。诸将谓齐愍王曰。﹁因而赴之。破燕必矣。﹂齐王因令人谓燕太子平曰。﹁寡人闻太子之义。将废私而立公。饬君臣之义。明父子之位。寡人之国小。不足以为先后。虽然。则唯太子所以令之。﹂太子因要党聚众。将军市被围公宫。攻子之。不克。将军市被及百姓反攻太子平。将军市被死。以徇。因构难数月。死者数万。众人恫恐。百姓离志。孟轲谓齐王曰。﹁今伐燕。此文。武之时。不可失也。﹂王因令章子将五都之兵。以因北地之众以伐燕。士卒不战。城门不闭。燕君哙死。齐大胜。燕子之亡二年。而燕人共立太子平。是为燕昭王。\\
\hspace*{2em}燕昭王于破燕之后即位。卑身厚币以招贤者。谓郭隗曰。﹁齐因孤之国乱而袭破燕。孤极知燕小力少。不足以报。然诚得贤士以共国。以雪先王之耻。孤之愿也。先生视可者。得身事之。﹂郭隗曰。﹁王必欲致士。先从隗始。况贤于隗者。岂远千里哉。﹂于是昭王为隗改筑宫而师事之。乐毅自魏往。邹衍自齐往。剧辛自赵往。士争趋燕。燕王吊死问孤。与百姓同甘苦。\\
\hspace*{2em}二十八年。燕国殷富。士卒乐轶轻战。于是遂以乐毅为上将军。与秦。楚。三晋合谋以伐齐。齐兵败。愍王出亡于外。燕兵独追北。入至临淄。尽取齐宝。烧其宫室宗庙。齐城之不下者。独唯聊。莒。即墨。其余皆属燕。六岁。\\
\hspace*{2em}昭王三十三年卒。子惠王立。\\
\hspace*{2em}惠王为太子时。与乐毅有隙。及即位。疑毅。使骑劫代将。乐毅亡走赵。齐田单以即墨击败燕军。骑劫死。燕兵引归。齐悉复得其故城。愍王死于莒。乃立其子为襄王。\\
\hspace*{2em}惠王七年卒。韩。魏。楚共伐燕。燕武成王立。\\
\hspace*{2em}武成王七年。齐田单伐我。拔中阳。十三年。秦败赵于长平四十余万。十四年。武成王卒。子孝王立。\\
\hspace*{2em}孝王元年。秦围邯郸者解去。三年卒。子今王喜立。\\
\hspace*{2em}今王喜四年。秦昭王卒。燕王命相栗腹约欢赵。以五百金为赵王酒。还报燕王曰。﹁赵王壮者皆死长平。其孤未壮。可伐也。﹂王召昌国君乐闲问之。对曰。﹁赵四战之国。其民习兵。不可伐。﹂王曰。﹁吾以五而伐一。﹂对曰。﹁不可。﹂燕王怒。群臣皆以为可。卒起二军。车二千乘。栗腹将而攻鄗。卿秦攻代。唯独大夫将渠谓燕王曰。﹁与人通关约交。以五百金饮人之王。使者报而反攻之。不祥。兵无成功。﹂燕王不听。自将偏军随之。将渠引燕王绶止之曰。﹁王必无自往。往无成功。﹂王蹴之以足。将渠泣曰。﹁臣非以自为。为王也。﹂燕军至宋子。赵使廉颇将。击破栗腹于鄗。乐乘破卿秦于代。乐闲奔赵。廉颇逐之五百余里。围其国。燕人请和。赵人不许。必令将渠处和。燕相将渠以处和。赵听将渠。解燕围。\\
\hspace*{2em}六年。秦灭东周。置三川郡。七年。秦拔赵榆次三十七城。秦置大原郡。九年。秦王政初即位。十年。赵使廉颇将攻繁阳。拔之。赵孝成王卒。悼襄王立。使乐乘代廉颇。廉颇不听。攻乐乘。乐乘走。廉颇奔大梁。十二年。赵使李牧攻燕。拔武遂。方城。剧辛故居赵。与庞暖善。已而亡走燕。燕见赵数困于秦。而廉颇去。令庞暖将也。欲因赵毙攻之。问剧辛。辛曰。﹁庞暖易与耳。﹂燕使剧辛将击赵。赵使庞暖击之。取燕军二万。杀剧辛。秦拔魏二十城。置东郡。十九年。秦拔赵之邺九城。赵悼襄王卒。二十三年。太子丹质于秦。亡归燕。二十五年。秦虏灭韩王安。置颍川郡。二十七年。秦虏赵王迁。灭赵。赵公子嘉自立为代王。\\
\hspace*{2em}燕见秦且灭六国。秦兵临易水。祸且至燕。太子丹阴养壮士二十人。使荆轲献督亢地图于秦。因袭刺秦王。秦王觉。杀轲。使将军王翦击燕。二十九年。秦攻拔我蓟。燕王亡。徙居辽东。斩丹以献秦。三十年。秦灭魏。\\
\hspace*{2em}三十三年。秦拔辽东。虏燕王喜。卒灭燕。是岁。秦将王贲亦虏代王嘉。\\
\hspace*{2em}太史公曰。召公奭可谓仁矣。甘棠且思之。况其人乎。燕外迫蛮貉。内措齐。晋。崎岖强国之间。最为弱小。几灭者数矣。然社稷血食者八九百岁。于姬姓独后亡。岂非召公之烈邪。\\

% ===== 索隐述赞（司马贞《史记索隐》）· 统一样式 =====
% 题头「索隐述赞曰」五字并入反白框；正文每四句一列，定宽盒 \hfil 居列中；
% \fontsize{24}{41} 列距对齐全书界行（121+41k）；\clearpage 使述赞独占一页。
\clearpage
\shuzanhead{索隐述赞曰}\linebreak
{\shuzanbody\fontsize{24}{41}\selectfont
\hbox to \hsize{\hfil 召伯作相。分陕而治。人惠其德。甘棠是思。\hfil}\linebreak
\hbox to \hsize{\hfil 庄送霸主。惠罗宠姬。文公从赵。苏秦骋辞。\hfil}\linebreak
\hbox to \hsize{\hfil 易王初立。齐宣我欺。燕哙无道。禅位子之。\hfil}\linebreak
\hbox to \hsize{\hfil 昭王待士。思报临淄。督亢不就。卒见芟夷。\hfil}%
\par}
\end{document}
